\chapter{Metodología}

Este capítulo describe cómo se planificó y ejecutó la investigación. Cubre las preguntas que motivaron el trabajo, el flujo de trabajo que estructuró la ejecución y los milestones que delimitaron las fases y sus entregables.

\section{Preguntas de investigación}

La investigación comenzó con preguntas exploratorias sobre el fenómeno, antes de tomar ninguna decisión de diseño. El punto de partida fue la observación de que las discusiones asíncronas en plataformas educativas no producen el tipo de aprendizaje que sí ocurre en un aula presencial, y que no estaba claro por qué ni qué se podría hacer al respecto.

\subsection{Fase exploratoria}

Las primeras preguntas apuntaban a entender el fenómeno:

\begin{itemize}
  \item ¿Qué ocurre en un aula presencial que no ocurre en una discusión asíncrona? ¿Por qué la dinámica social del aprendizaje no se traslada a los foros de texto?
  \item ¿Qué hace a un facilitador efectivo? ¿Qué acciones concretas produce un instructor cuando interviene en una discusión?
  \item ¿Cómo se ve el proceso de intervención desde dentro? ¿Qué información necesita un facilitador para tomar una decisión?
  \item ¿Es posible hacer ese proceso observable y controlable, o es opaco por naturaleza?
\end{itemize}

Estas preguntas guiaron la revisión de la literatura y la exploración técnica inicial. La respuesta a la cuarta pregunta, en particular, determinó la dirección del diseño, puesto que si el proceso de intervención puede describirse como un conjunto de condiciones observables más una acción correspondiente puede implementarse como un pipeline; si no puede, la automatización no es viable. La evidencia de los sistemas de tutoría inteligente \autocite{VanLehn2011} y la literatura sobre roles de facilitación \autocite{PilkingtonWalker2003} sugirió que sí es posible, dentro de límites.

\subsection{Fase de diseño}

Una vez comprendido el fenómeno, las preguntas se volvieron más concretas:

\begin{itemize}
  \item ¿Cómo pueden operacionalizarse los roles de facilitación como agentes LLM en un pipeline de análisis de discusiones?
  \item ¿Qué criterios permiten detectar automáticamente el estado de una discusión para determinar si y cuándo intervenir?
  \item ¿Qué criterios permiten evaluar si una intervención generada por IA es pedagógicamente adecuada?
  \item ¿Puede este sistema integrarse en una plataforma de educación abierta sin modificar la plataforma base?
\end{itemize}

Estas preguntas estructuran el capítulo de desarrollo, ya que cada una corresponde a un componente del sistema o a una decisión de diseño documentada.

\section{Flujo de trabajo}

El proyecto usa un flujo de trabajo de I+D inspirado en prácticas ágiles, con kanban como herramienta de seguimiento, pero sin sprints fijos. La razón es que la incertidumbre de investigación no se ajusta bien a iteraciones de longitud predeterminada, ya que algunas preguntas se resuelven en días mientras otras requieren semanas de exploración antes de que sea posible tomar una decisión.

En la práctica, el trabajo se organiza en torno a ciclos de descubrimiento, donde cada fase produce hallazgos que informan la siguiente. Las herramientas principales son GitHub issues, hitos y Architecture Decision Records (ADRs):

\begin{itemize}
  \item Los \textbf{issues} capturan lo que debe explorarse o implementarse, con contexto suficiente para retomar el trabajo después de un tiempo.
  \item Los \textbf{ADRs} registran cada decisión de diseño significativa, su razonamiento y las alternativas descartadas. Son la memoria del proyecto.
  \item Los \textbf{hitos} delimitan fases. Una fase cierra cuando las preguntas abiertas que la definen tienen respuesta, no cuando expira un plazo.
\end{itemize}

En esta memoria solo se reportan decisiones de diseño respaldadas por evidencia trazable en resultados, tablas, figuras o ADRs del proyecto.

Antes de comprometerse con una implementación, abordamos cada área de incertidumbre técnica mediante un spike o prueba de concepto. Los spikes no tienen calidad de producción; su objetivo es reducir incertidumbre y documentar hallazgos, no entregar código mantenible.

\subsection{Uso de asistentes de IA en la implementación}

La implementación técnica la realicé con apoyo de asistentes de programación basados en LLM, en particular Claude y Codex, como herramientas de productividad para explorar alternativas de código, acelerar prototipos y revisar consistencia entre componentes. Su uso siguió un ciclo iterativo de \textit{propuesta, revisión y corrección}: el asistente proponía cambios o borradores, yo evaluaba críticamente su validez respecto a los objetivos pedagógicos y de arquitectura, y después refinaba la solución en nuevas iteraciones.

Este esquema no delegó las decisiones de diseño del artefacto. La definición de requisitos, la elección de arquitectura, los criterios de evaluación y la aceptación final de cada cambio quedaron bajo mi responsabilidad directa. En términos metodológicos, los asistentes funcionaron como soporte de implementación y análisis, no como fuente autónoma de decisiones de investigación.

\section{Planificación por hitos}

El proyecto se estructura en seis hitos que corresponden a las fases de un proyecto de I+D aplicado:

\begin{table}[h]
\centering
\begin{tabular}{lp{9cm}}
\toprule
\textbf{Hito} & \textbf{Descripción} \\
\midrule
1. Definición del problema & Formalización del problema, objetivos, requisitos y enfoque metodológico. \\
2. Investigación y exploración & Pruebas de concepto para reducir incertidumbre técnica antes del diseño. \\
3. Implementación & Diseño e implementación del sistema a partir de los hallazgos del hito anterior. \\
4. Integración & Integración del servicio de facilitación con Open edX como prueba de concepto en contexto real. \\
5. Control de calidad y feedback & Validación, recopilación de feedback y preparación de la documentación final. \\
6. Despliegue e infraestructura & Puesta en operación en dos entornos (experimentos y ciclo live), incluyendo decisiones de arquitectura de despliegue y operación. \\
\bottomrule
\end{tabular}
\caption{Hitos del proyecto y su descripción.}
\label{tab:milestones}
\end{table}

Los hitos 1 y 2 corresponden a la fase exploratoria de la investigación. Los hitos 3 y 4 corresponden a la fase de diseño e implementación. El hito 5 corresponde a la evaluación y el hito 6 a la operacionalización del artefacto en entornos reales de ejecución. Decidimos esta extensión de cinco a seis hitos durante la ejecución del proyecto, al comprobar que las decisiones de despliegue condicionaban directamente la validez de la prueba de concepto de integración. El capítulo de desarrollo sigue esta estructura final.
